\documentclass[12pt]{article}
\usepackage[T1]{fontenc}
\usepackage[utf8]{inputenc}
\usepackage{lmodern}
\usepackage{textcomp}
\usepackage{enumitem}
\usepackage{geometry}
\geometry{margin=1in}
\begin{document}
\begin{center}
Answer Key - Business Administration - Graduate Level Assessment 16 \
\small (Answers are ordered exactly as in the question paper.)
\end{center}

\section*{Multiple Choice}
\begin{enumerate}[label=\arabic*.]
\item \textbf{Answer:} A
\\ \textit{Options:} A) \$1.11; B) \$1.50; C) \$1.93; D) \$0.58; E) \$0.96
\\ \textit{Note:} The correct answer is $1.11 because using the 6%, 60-day method, the formula for simple interest (I = P × r × t) applies with P as $193, r as 0.06, and t as 38/60, resulting in the calculation of interest as \$1.11.
\item \textbf{Answer:} A
\\ \textit{Options:} A) Task-oriented; B) Servant; C) Transactional; D) Laissez-faire; E) Bureaucratic
\\ \textit{Note:} The correct answer is "task-oriented" because this leadership style emphasizes the importance of considering team members' views and ideas to effectively achieve group goals and enhance productivity.
\item \textbf{Answer:} D
\\ \textit{Options:} A) Customer participation and environmental acquisition.; B) Environmental acquisition and environmental relationship.; C) Customer retention and environmental relationship.; D) Customer participation and environmental relationship.; E) Customer acquisition and customer retention.
\\ \textit{Note:} The correct answer is based on Pine and Gilmore's framework, which identifies customer participation and environmental relationship as the two dimensions that define the four realms of experience: entertainment, education, escapism, and esthetic.
\item \textbf{Answer:} A
\\ \textit{Options:} A) equitably manage the sales bonus pool; B) increase the competitiveness within the sales team; C) create multiple points of contact between the customer and the company; D) distribute the work involved in managing an account; E) reduce the workload of individual salespeople
\\ \textit{Note:} The primary objective of team-based selling is to ensure that sales incentives are fairly distributed among team members, thereby promoting collaboration and alignment towards common sales goals.
\item \textbf{Answer:} E
\\ \textit{Options:} A) \$56; B) \$26; C) \$16; D) \$18; E) \$24
\\ \textit{Note:} The net price is calculated by subtracting the trade discount of $16 from the list price of $40, resulting in a net price of \$24.
\end{enumerate}

\section*{True / False}
\begin{enumerate}[label=\arabic*.]
\setcounter{enumi}{5}
\item \textbf{Answer:} True
\\ \textit{Options:} A) True; B) False
\\ \textit{Note:} The percentage discount is calculated by taking the difference between the list price and the net price, which is ($300 - $174) / \$300 = 0.42 or 42\%.
\item \textbf{Answer:} True
\\ \textit{Options:} A) True; B) False
\\ \textit{Note:} The statement is true because a trade discount is indeed calculated by subtracting the net price from the list price, and in this case, that difference is accurately stated as \$3.95.
\item \textbf{Answer:} True
\\ \textit{Options:} A) True; B) False
\\ \textit{Note:} The statement is true because the conversions of 678 cm to approximately 22 feet, 316 liters to about 575.12 pints, and 50 kg to 110 pounds are accurate based on standard measurement conversion factors.
\item \textbf{Answer:} False
\\ \textit{Options:} A) True; B) False
\\ \textit{Note:} The statement is false because the cost difference between a 3-year policy and three individual one-year policies typically reflects factors such as premium discounts for multi-year commitments, which would likely result in a difference greater than \$63.00.
\item \textbf{Answer:} True
\\ \textit{Options:} A) True; B) False
\\ \textit{Note:} The statement is true because the rate of return can be calculated using the formula for the annualized return on investment, which shows that repaying $800 on a $500 loan after five years results in an effective annual return of approximately 13.86\%.
\end{enumerate}

\section*{Long Form Response}
\begin{enumerate}[label=\arabic*.]
\setcounter{enumi}{10}
\item \textbf{Answer:} An$_{80}$\% coinsurance clause in property insurance requires that the property be insured for at least 80\% of its value to receive full compensation for a loss. In this scenario, the house is valued at $10,000 but is only insured for $7,000, which does not meet the 80\% requirement (i.e., $8,000). Consequently, in the event of an $8,500 loss, the payout will be reduced due to the underinsurance. The calculation for the payout is determined by the formula: (Amount Insured / Required Amount Insured) x Loss Amount, resulting in ($7,000 / $8,000) x $8,500 = $6,125. Therefore, the insurance payout is $6,125, leading to an owner's loss of $2,375 when considering the total loss against the payout. This situation underscores the importance of adequately insuring property to avoid significant financial penalties when claims are made.
\\ \textit{Note:} correct because it accurately identifies that the property must be insured for at least 80\% of its value to qualify for full compensation, and since the house is underinsured, it applies the coinsurance penalty to determine the reduced payout based on the calculated ratio.
\item \textbf{Answer:} Lateral communication plays a crucial role in organizations by facilitating the exchange of information and ideas among employees at the same hierarchical level. During official meetings, this type of communication becomes particularly significant as it provides a structured environment for employees to share insights, discuss challenges, and collaborate on solutions. Although information sharing can occur informally, the structured nature of meetings ensures that critical updates and collaborative discussions happen systematically. This intentional sharing not only enhances teamwork but also fosters a culture of transparency and trust within the organization. Ultimately, effective lateral communication during meetings directly contributes to improved collaboration and the efficient flow of information among employees, leading to better organizational performance.
\\ \textit{Note:} correct because it emphasizes that lateral communication fosters collaboration and information sharing among employees at the same hierarchical level, particularly during official meetings, which provide a formal setting for structured dialogue and problem-solving.
\end{enumerate}

\vfill
\noindent\textit{End of Answer Key}
\end{document}